\chapter{Construction de portefeuille}\label{ch:iv2:portfolio-construction}

\begin{ivtrains}
Construire un portefeuille à la main, du cas à deux actifs jusqu'aux contraintes, à la parité de risque, aux modèles factoriels, à la loi fondamentale, au shrinkage et au turnover. Les rendements sont annuels et en excès de cash, sauf indication contraire dans un énoncé, et tous les nombres sont inventés.
\end{ivtrains}

\begin{question}[$\star$ \iqroles{QR, QT, Analyste} \iqfirm{toutes} \iqlevel{junior}]\label{iq:iv2:portfolio-construction:1}
Deux actions ont des volatilités de 20~\% et 30~\%, et une corrélation de 0.5. Quelles pondérations donnent le portefeuille de variance minimale~? Quelle est sa volatilité~?
\end{question}

\begin{question}[$\star$ \iqroles{QR, QT, Analyste} \iqfirm{toutes} \iqlevel{junior}]\label{iq:iv2:portfolio-construction:2}
Vous disposez de deux actifs de même volatilité, 20~\%, et d'une corrélation de 0.2. Quel est le mélange de variance minimale, et quelle est sa volatilité~?
\end{question}

\begin{question}[$\star$ \iqroles{QR, Risque, Analyste} \iqfirm{toutes} \iqlevel{junior}]\label{iq:iv2:portfolio-construction:3}
Deux actifs ont des volatilités de 10~\% et 20~\%. Quelles pondérations rendent égales leurs contributions au risque~? La réponse dépend-elle de la corrélation~?
\end{question}

\begin{question}[$\star$ \iqroles{QR, QT} \iqfirm{hedge fund} \iqlevel{junior}]\label{iq:iv2:portfolio-construction:4}
Un signal a un coefficient d'information de 0.05. En utilisant la loi fondamentale, quel est le ratio d'information avec 100 paris indépendants par an~? Avec 252~?
\end{question}

\begin{question}[$\star$ \iqroles{QR, QT, Trader} \iqfirm{hedge fund, teneur de marché} \iqlevel{junior}]\label{iq:iv2:portfolio-construction:5}
Un portefeuille passe de pondérations (0.5, 0.3, 0.2) à (0.4, 0.4, 0.2). Quel est le turnover à sens unique, et que payez-vous à 10 points de base par dollar échangé~? Si vous faites cela chaque mois, quel en est le coût annuel~?
\end{question}

\begin{question}[$\star$ \iqroles{QT, Trader, Risque} \iqfirm{hedge fund} \iqlevel{junior}]\label{iq:iv2:portfolio-construction:6}
Vous êtes acheteur de 10 millions de dollars d'une action de bêta 1.2. Vous vous couvrez avec une action de bêta 0.8. Quel montant vendez-vous à découvert pour une neutralité en bêta~? Quelle est votre position nette en dollars~?
\end{question}

\begin{question}[$\star\star$ \iqroles{QR, QT} \iqfirm{hedge fund, banque} \iqlevel{junior, confirmé}]\label{iq:iv2:portfolio-construction:7}
Deux actifs ont des excès de rendement espérés de 8~\% et 10~\%, des variances de 0.04 et 0.09, et une covariance de 0.03. Quelles sont les pondérations du portefeuille tangent et son ratio de Sharpe~? Pour un investisseur moyenne--variance d'aversion au risque 4, que détient-il, et quels sont son rendement et sa volatilité~?
\end{question}

\begin{question}[$\star\star$ \iqroles{QR, QT, Risque} \iqfirm{toutes} \iqlevel{junior, confirmé}]\label{iq:iv2:portfolio-construction:8}
Deux actions ont des volatilités de 20~\% et 30~\%, et une corrélation de 0.8. Quel est le portefeuille de variance minimale~? Pourquoi une des pondérations est-elle négative~? Donnez la condition exacte.
\end{question}

\begin{question}[$\star\star$ \iqroles{QR, QT} \iqfirm{hedge fund} \iqlevel{junior, confirmé}]\label{iq:iv2:portfolio-construction:9}
L'actif A a un excès de rendement de 10~\% et une volatilité de 20~\%. L'actif B a 5~\% et 30~\%. Leur covariance est de 0.05. Calculez le ratio de Sharpe du portefeuille tangent sans contrainte, et donnez ses pondérations sous la forme $\Sigma^{-1}\mu$, soit les positions pour une aversion au risque de 1. Quel est le meilleur ratio de Sharpe en positions longues uniquement, et comment le savez-vous~?
\end{question}

\begin{question}[$\star\star$ \iqroles{QR, Risque, Analyste} \iqfirm{toutes} \iqlevel{junior}]\label{iq:iv2:portfolio-construction:10}
Des actions ont chacune une volatilité de 30~\% et une corrélation deux à deux commune de 0.3. Quelle est la volatilité d'un portefeuille à pondérations égales de 10, 30 et 100 de ces actions~? Quelle est la limite~? Combien d'actions faut-il pour être à moins de 1 point de cette limite~?
\end{question}

\begin{question}[$\star\star$ \iqroles{QR, QT} \iqfirm{hedge fund} \iqlevel{junior, confirmé}]\label{iq:iv2:portfolio-construction:11}
Vous classez 500 actions chaque année avec un coefficient d'information de 0.02, et les erreurs de prévision ont une corrélation deux à deux de 0.01. Quelle sont la largeur effective et le ratio d'information~? Comparez avec ce que donne la loi fondamentale si on les traite comme indépendantes.
\end{question}

\begin{question}[$\star\star$ \iqroles{QR, QT} \iqfirm{hedge fund} \iqlevel{junior, confirmé}]\label{iq:iv2:portfolio-construction:12}
Un signal a un coefficient d'information de 0.05 et 500 paris indépendants. Le coefficient de transfert du portefeuille contraint en positions longues uniquement vaut 0.3. Quel est le ratio d'information espéré~? Que mesure ce coefficient, et comment l'amélioreriez-vous~?
\end{question}

\begin{question}[$\star\star$ \iqroles{QR, QT, Analyste} \iqfirm{hedge fund, banque} \iqlevel{junior, confirmé}]\label{iq:iv2:portfolio-construction:13}
Deux actifs ont des moyennes empiriques d'excès de rendement de 15~\% et $-5~\%$, sur 20 ans de données, avec une volatilité de 20~\% chacun. Quelle est l'erreur type de chaque moyenne~? Contractez les deux moitiés vers un 5~\% commun. Que se passe-t-il~? Pourquoi l'optimisation moyenne--variance aggrave-t-elle l'erreur d'estimation~?
\end{question}

\begin{question}[$\star\star$ \iqroles{QR, Risque} \iqfirm{hedge fund, banque} \iqlevel{junior, confirmé}]\label{iq:iv2:portfolio-construction:14}
Une covariance empirique a pour valeurs propres 5, 1 et 0.01. Quel est son conditionnement~? Contractez-la de 20~\% vers un multiple de l'identité de même trace. Quelles sont les nouvelles valeurs propres et le nouveau conditionnement~? Pourquoi cela aide-t-il un optimiseur de variance minimale~?
\end{question}

\begin{question}[$\star\star$ \iqroles{QR, QT} \iqfirm{hedge fund} \iqlevel{junior, confirmé}]\label{iq:iv2:portfolio-construction:15}
Vous calculez une matrice de corrélation pour 100 actions à partir de 1\,000 jours de pur bruit. Où se situent les valeurs propres, selon Marchenko-Pastur~? Qu'arrive-t-il avec 500 actions~? Comment en tirez-vous parti~?
\end{question}

\begin{question}[$\star\star$ \iqroles{QR, Risque, QD} \iqfirm{hedge fund, banque} \iqlevel{junior, confirmé}]\label{iq:iv2:portfolio-construction:16}
Pour 500 actions, combien de paramètres comporte une matrice de covariance complète~? Combien en comporte un modèle à 10 facteurs, avec une covariance complète entre facteurs et des variances idiosyncratiques individuelles~? Quelle hypothèse permet cela, et quand échoue-t-elle~?
\end{question}

\begin{question}[$\star\star$ \iqroles{QR, QT, Analyste} \iqfirm{toutes} \iqlevel{junior, confirmé}]\label{iq:iv2:portfolio-construction:17}
Deux actifs indépendants ont chacun un rendement arithmétique espéré nul et une volatilité de 20~\%, en loi lognormale. Un portefeuille est rééquilibré chaque année à 50/50. Quelle est sa croissance logarithmique espérée, comparée à celle de chaque actif~? Comment se comparent une stratégie de rééquilibrage et une stratégie de conservation sur 30 ans~? Qu'en est-il si la corrélation vaut 1~?
\end{question}

\begin{question}[$\star\star\star$ \iqroles{QR, QT} \iqfirm{hedge fund} \iqlevel{confirmé}]\label{iq:iv2:portfolio-construction:18}
Dix actifs ont chacun une moyenne vraie de 0.5~\% par mois, une volatilité de 5~\%, et une corrélation deux à deux de 0.3. Vous estimez le portefeuille tangent à partir de 60 mois de données. Quel est le Sharpe vrai des pondérations égales~? Quels Sharpe annualisés, dans l'échantillon et hors échantillon, obtient-on à peu près avec les pondérations estimées~? Que retenez-vous de cela~?
\end{question}

\begin{question}[$\star\star\star$ \iqroles{QR, QT, Analyste} \iqfirm{hedge fund, banque} \iqlevel{confirmé}]\label{iq:iv2:portfolio-construction:19}
Les pondérations de marché sont de 0.6 et 0.4, les variances de 0.04 et 0.09, la covariance de 0.03 et l'aversion au risque de 2.5. Quels sont les excès de rendement implicites~? Vous avez l'opinion que l'actif 1 rapporte 10~\%, avec $\tau = 0.05$ et une incertitude sur cette opinion égale à $\tau$ fois sa variance. Donnez les rendements et les pondérations a posteriori de Black-Litterman. Que se passe-t-il lorsque l'incertitude tend vers zéro~?
\end{question}

\begin{question}[$\star\star\star$ \iqroles{QR, Risque} \iqfirm{hedge fund, banque} \iqlevel{junior, confirmé}]\label{iq:iv2:portfolio-construction:20}
Trois actifs ont des variances de 0.04, 0.01 et 0.09. La covariance vaut 0.01 entre les deux premiers, 0.004 entre les deux derniers, et 0 entre le premier et le troisième. Calculez les parts de contribution au risque pour des pondérations égales, pour des pondérations inverses à la volatilité, c'est-à-dire proportionnelles à un sur la volatilité de chaque actif, et pour (0.2763, 0.5119, 0.2117). Lequel est le portefeuille à risque égal~? Que possède-t-il que la volatilité inverse n'a pas~?
\end{question}

\begin{question}[$\star\star\star$ \iqroles{QR, QD} \iqfirm{toutes} \iqlevel{confirmé}]\label{iq:iv2:portfolio-construction:21}
Trois actifs ont des excès de rendement espérés de 6~\%, 8~\% et 12~\%, et des variances de 0.04, 0.09 et 0.0225. Leurs covariances valent 0.012, 0.012 et 0.018, pour les couples (1, 2), (1, 3) et (2, 3). Calculez le rendement et la volatilité du portefeuille de variance minimale globale. Trouvez le portefeuille de la frontière dont le rendement vaut 12~\%. Comment le théorème des deux fonds se manifeste-t-il ici~?
\end{question}

\begin{question}[$\star\star\star$ \iqroles{QR, QD, Risque} \iqfirm{toutes} \iqlevel{confirmé}]\label{iq:iv2:portfolio-construction:22}
Trois actifs ont des variances de 0.04, 0.09 et 0.0225, et des covariances de 0.012, 0.012 et 0.018 pour les couples (1, 2), (1, 3) et (2, 3). Avec des pondérations de somme égale à 1, leurs pondérations de variance minimale sans contrainte valent (0.268, 0.043, 0.689). Imposez un poids maximal de 40~\%. Quelle est la solution~? Comment savez-vous qu'elle est optimale, par les conditions et non par un solveur~?
\end{question}

\begin{question}[$\star\star\star$ \iqroles{QR, QT, Trader} \iqfirm{hedge fund, teneur de marché} \iqlevel{junior, confirmé}]\label{iq:iv2:portfolio-construction:23}
Un signal indique un rendement espéré de 2 points de base par jour, qui décroît avec une demi-vie de 5 jours. Les coûts de transaction valent 5 points de base dans chaque sens. L'opération en vaut-elle la peine~? Quel est le signal initial de seuil de rentabilité~? Que change-t-il si la position doit être vendue au bout d'un jour~?
\end{question}

\begin{question}[$\star\star\star$ \iqroles{QR, QT} \iqfirm{hedge fund} \iqlevel{junior, confirmé}]\label{iq:iv2:portfolio-construction:24}
Deux stratégies ont des ratios de Sharpe de 1.0 et 0.8, et une corrélation de 0.3. Quel est le Sharpe du meilleur mélange, et avec quelle allocation de risque~? À quelle corrélation la seconde stratégie n'apporte-t-elle rien~? Que se passe-t-il pour une corrélation négative~?
\end{question}

\begin{question}[$\star\star\star\star$ \iqroles{QR, QT, Risque} \iqfirm{hedge fund} \iqlevel{confirmé}]\label{iq:iv2:portfolio-construction:25}
Vous disposez de trois actions avec des rendements espérés de 6~\%, 3~\% et $-4~\%$, des volatilités de 25~\%, 20~\% et 30~\%, une corrélation commune de 0.4, et des bêtas de 1.2, 0.8 et 1.0. Votre aversion au risque vaut 5. Trouvez le portefeuille sans contrainte et son bêta. Rendez-le ensuite neutre en bêta, puis aussi neutre en dollars. Comparez rendement, volatilité et Sharpe. Ramenez-le à une volatilité de 10~\%. Ajoutez ensuite un coût de 0.2~\% par dollar échangé à partir d'une position nulle. Que change-t-il, et comment vérifiez-vous la réponse sans solveur~?
\end{question}
